% ==============================================================================
% ภาคผนวก: รายการหลักฐานประกอบการประเมิน
% ==============================================================================

\partpage{ส่วนที่ 4}{ภาคผนวก}

\section{รายการหลักฐานประกอบการประเมิน}

รายงานการประเมินตนเองฉบับนี้จัดทำขึ้นเพื่อเป็นหลักฐานประกอบการประกันคุณภาพการศึกษาภายใน ระดับหลักสูตร ตามเกณฑ์ AUN-QA Version 4 ของหลักสูตรวิศวกรรมศาสตรมหาบัณฑิต สาขาวิชาวิศวกรรมคอมพิวเตอร์และปัญญาประดิษฐ์ มหาวิทยาลัยราชภัฏอุตรดิตถ์ ประจำปีการศึกษา 2568 รายละเอียดของหลักฐานที่อ้างอิงในรายงานฉบับนี้จัดแสดงไว้ในตารางต่อไปนี้

\setlength{\tabcolsep}{3pt}
\begin{xltabular}{\linewidth}{|>{\centering\arraybackslash}p{1.1cm}|>{\centering\arraybackslash}p{2.0cm}|Y|>{\centering\arraybackslash}p{2.0cm}|Y|}
\caption{รายการหลักฐานประกอบการประเมิน (Evidence List)} \\
\hline
\textbf{ลำดับ} & \textbf{รหัส} & \textbf{ชื่อหลักฐาน} & \textbf{สถานะ} & \textbf{หมายเหตุ} \\
\hline
\endfirsthead
\hline
\textbf{ลำดับ} & \textbf{รหัส} & \textbf{ชื่อหลักฐาน} & \textbf{สถานะ} & \textbf{หมายเหตุ} \\
\hline
\endhead
\hline
\endfoot
1 & AUNQA-1-1 & เล่ม มคอ.2 หลักสูตรวิศวกรรมคอมพิวเตอร์และปัญญาประดิษฐ์ ฉบับผ่านสภา & \cellcolor{green!25}มี & หลักสูตร wiki \url{https://uru.ac.th/curriculum/master-eng-ai} \\
2 & AUNQA-1-2 & รายงานการประชุมคณะกรรมการประจำหลักสูตร & \cellcolor{green!25}มี & --- \\
3 & AUNQA-1-3 & รวมเล่ม CLO รายวิชา พร้อม CLO--PLO Mapping & \cellcolor{green!25}มี & หลักสูตร wiki \url{https://uru.ac.th/curriculum/master-eng-ai} \\
4 & AUNQA-1-4 & มติสภามหาวิทยาลัยราชภัฏอุตรดิตถ์ ครั้งที่ 4/2568 & \cellcolor{green!25}มี & มติ 4 เมษายน 2568 \\
5 & AUNQA-2-1 & มคอ.2 ฉบับผ่านสภามหาวิทยาลัยราชภัฏอุตรดิตถ์ & \cellcolor{green!25}มี & หลักสูตร wiki \url{https://uru.ac.th/curriculum/master-eng-ai} \\
6 & AUNQA-2-2 & รวมเล่ม CLO รายวิชา พร้อม CLO--PLO Mapping & \cellcolor{green!25}มี & หลักสูตร wiki \url{https://uru.ac.th/curriculum/master-eng-ai} \\
7 & AUNQA-2-3 & ข้อบังคับมหาวิทยาลัยราชภัฏอุตรดิตถ์ ว่าด้วยการจัดการศึกษาระดับบัณฑิตศึกษา พ.ศ. 2566 & \cellcolor{green!25}มี & ฝ่ายวิชาการ มร.อ.\ \url{https://academic.uru.ac.th/document/02898.pdf} \\
8 & AUNQA-2-4 & มคอ.3 ของรายวิชาบังคับ (อยู่ระหว่างจัดทำ) & \cellcolor{red!25}MARK & กำลังจัดทำ \\
9 & AUNQA-3-1 & มคอ.2 หมวด 4 (ลักษณะและการดำเนินการ) และหมวด 5 (แผนการสอนและการประเมิน) & \cellcolor{green!25}มี & หลักสูตร wiki \url{https://uru.ac.th/curriculum/master-eng-ai} \\
10 & AUNQA-3-2 & รวมเล่ม CLO รายวิชา — ตาราง Learning Environment ครบ 32 วิชา & \cellcolor{green!25}มี & หลักสูตร wiki \url{https://uru.ac.th/curriculum/master-eng-ai} \\
11 & AUNQA-3-3 & มคอ.3 ของรายวิชาบังคับ (อยู่ระหว่างจัดทำ) & \cellcolor{red!25}MARK & กำลังจัดทำ \\
12 & AUNQA-3-4 & ข้อบังคับมหาวิทยาลัยฯ ว่าด้วยการจัดการศึกษาระดับบัณฑิตศึกษา พ.ศ. 2566 & \cellcolor{green!25}มี & ฝ่ายวิชาการ มร.อ.\ \url{https://academic.uru.ac.th/document/02898.pdf} \\
13 & AUNQA-3-5 & แผนการสอนรายวิชาบังคับ 6 วิชา (Constructive Alignment แต่ละ CLO) & \cellcolor{red!25}MARK & ต้องจัดทำ \\
14 & AUNQA-4-1 & ประกาศมหาวิทยาลัยฯ เรื่อง ระบบและแนวทางการอุทธรณ์ผลการประเมินฯ พ.ศ.~2567 & \cellcolor{green!25}มี & \url{https://academic.uru.ac.th/document/06943.pdf} \\
15 & AUNQA-4-2 & ข้อบังคับมหาวิทยาลัยฯ ว่าด้วยการจัดการศึกษาระดับบัณฑิตศึกษา พ.ศ.~2566 & \cellcolor{green!25}มี & ฉบับหลัก \url{https://academic.uru.ac.th/document/02898.pdf}; ฉบับที่ 2 \url{https://academic.uru.ac.th/document/02897.pdf} \\
16 & AUNQA-4-3 & ประกาศมหาวิทยาลัยฯ เรื่อง แนวปฏิบัติและการกำกับติดตามการส่งผลการประเมินผลการศึกษา พ.ศ.~2567 & \cellcolor{green!25}มี & \url{https://academic.uru.ac.th/document/06941.pdf} \\
17 & AUNQA-5-1 & แผนบริหารทรัพยากรบุคคล มหาวิทยาลัยราชภัฏอุตรดิตถ์ พ.ศ.~2565--2569 & \cellcolor{yellow!25}มีบางส่วน & \url{https://personnel.uru.ac.th/wp-content/uploads/2022/03/plan_manage_2565-2569.pdf} \\
18 & AUNQA-5-2 & ประกาศมหาวิทยาลัยฯ เรื่อง นโยบายและแนวทางการพัฒนาคุณภาพอาจารย์ & \cellcolor{yellow!25}มีบางส่วน & \url{https://personnel.uru.ac.th/wp-content/uploads/2022/03/plan_hr_2565-2569.pdf} \\
19 & AUNQA-5-3 & ข้อบังคับมหาวิทยาลัยฯ ว่าด้วยการบริหารงานบุคคลของพนักงานมหาวิทยาลัย & \cellcolor{green!25}มี & งานบุคคล มร.อ.\ \url{https://personnel.uru.ac.th/} \\
20 & AUNQA-5-4 & รายงานการพัฒนาอาจารย์/บันทึกการเข้าร่วมอบรมของอาจารย์ผู้รับผิดชอบหลักสูตร & \cellcolor{yellow!25}มีบางส่วน & ต้องรวบรวมให้ครบ \\
21 & AUNQA-6-1 & ข้อบังคับมหาวิทยาลัยฯ ว่าด้วยการจัดการศึกษาระดับบัณฑิตศึกษา พ.ศ.~2566 & \cellcolor{green!25}มี & ฝ่ายวิชาการ มร.อ.\ \url{https://academic.uru.ac.th/document/02898.pdf} \\
22 & AUNQA-6-2 & ประกาศการรับสมัครนักศึกษาบัณฑิตศึกษา ปีการศึกษา 2568 & \cellcolor{yellow!25}มีบางส่วน & ตามรอบบัณฑิตวิทยาลัย \url{https://academic.uru.ac.th/SmartUru/} \\
23 & AUNQA-6-3 & รายชื่อฐานข้อมูลและบริการของสำนักวิทยบริการ ม.ราชภัฏอุตรดิตถ์ & \cellcolor{green!25}มี & สำนักวิทยบริการ \url{https://sites.google.com/uru.ac.th/online-database} \\
24 & AUNQA-6-4 & รายงานความพึงพอใจของนักศึกษาบัณฑิตศึกษา / แบบสำรวจเฉพาะหลักสูตร & \cellcolor{red!25}MARK & อยู่ระหว่างพัฒนาเครื่องมือ \\
25 & AUNQA-6-5 & ปฏิทินกิจกรรมเสริมและหลักฐานการจัดกิจกรรม & \cellcolor{yellow!25}มีบางส่วน & เติมเมื่อเริ่มจัดกิจกรรมจริง \\
26 & AUNQA-7-1 & รายงานสิ่งอำนวยความสะดวกของคณะเทคโนโลยีอุตสาหกรรม (อาคาร 10) & \cellcolor{yellow!25}มีบางส่วน & ต้องยืนยันความพร้อม \\
27 & AUNQA-7-2 & รายชื่อฐานข้อมูลดิจิทัล/สัญญาใช้งานของสำนักวิทยบริการ & \cellcolor{green!25}มี & ห้องสมุด/OPAC \url{https://opac.uru.ac.th}; ฐานข้อมูล \url{https://sites.google.com/uru.ac.th/online-database} \\
28 & AUNQA-7-3 & แผนแม่บทเทคโนโลยีสารสนเทศของศูนย์คอมพิวเตอร์ & \cellcolor{yellow!25}มีบางส่วน & แผนปฏิบัติการสำนักวิทยบริการฯ ปี 2568 \url{https://arit.uru.ac.th/upload/plan/Plan2568.pdf} \\
29 & AUNQA-7-4 & รายงานผลประเมินคุณภาพสิ่งอำนวยความสะดวก/บริการของสำนักวิทยบริการ/ศูนย์คอมพิวเตอร์ & \cellcolor{yellow!25}มีบางส่วน & --- \\
30 & AUNQA-7-5 & บันทึกการตรวจสอบความพร้อมห้องปฏิบัติการ/อุปกรณ์ก่อนเปิดภาค & \cellcolor{yellow!25}มีบางส่วน & เติมเมื่อมีข้อมูลจริง \\
31 & AUNQA-8-1 & ระบบรายงานสถิตินักศึกษาของกองบริการการศึกษา (ข้อมูล REG) & \cellcolor{yellow!25}มีบางส่วน & นักศึกษา 4 คน; ระบบลงทะเบียน \url{https://bundit.uru.ac.th/} \\
32 & AUNQA-8-2 & แผน Alumni Tracking System ของหลักสูตร CE\&AI & \cellcolor{red!25}MARK & ต้องจัดทำ \\
33 & AUNQA-8-3 & รายงานผลงานวิจัยของอาจารย์ผู้รับผิดชอบหลักสูตรย้อนหลัง 5 ปี & \cellcolor{yellow!25}มีบางส่วน & อยู่ระหว่างรวบรวม \\
34 & AUNQA-8-4 & สรุปผล CLO / ความก้าวหน้าวิทยานิพนธ์และข้อมูลระหว่างทางเชื่อม PLO & \cellcolor{yellow!25}มีบางส่วน & สะสมตามรอบรายวิชา \\
35 & AUNQA-8-5 & แบบสำรวจความพึงพอใจผู้มีส่วนได้ส่วนเสีย (แผนและผลรายภาค) & \cellcolor{yellow!25}มีบางส่วน & รอสรุปผลครบ \\
36 & AUNQA-9-1 & รายงานการประชุมคณะกรรมการหลักสูตร & \cellcolor{green!25}มี & มีรายละเอียดการพัฒนาหลักสูตร \\
37 & AUNQA-9-2 & ข้อมูล Vision/Mission มหาวิทยาลัยและคณะ & \cellcolor{green!25}มี & เว็บไซต์ มหาวิทยาลัย \url{https://www.uru.ac.th}; คณะเทคโนโลยีอุตสาหกรรม \url{https://industrial.uru.ac.th} \\
38 & AUNQA-9-3 & ยุทธศาสตร์ 5 ปีของมหาวิทยาลัยและคณะ & \cellcolor{green!25}มี & ฝ่ายแผนและความร่วมมือ \url{https://plan.uru.ac.th} \\
39 & AUNQA-9-4 & ข้อมูลจำนวนนักศึกษาและบัณฑิต & \cellcolor{yellow!25}มีบางส่วน & นักศึกษา 4 คน (2568) ยังไม่มีบัณฑิต \\
\hline
\end{xltabular}

\bigskip

\section{ความหมายของสถานะหลักฐาน}

สถานะหลักฐานในรายงานนี้แบ่งออกเป็น 3 ระดับ ได้แก่ \textbf{มี} หมายถึงหลักฐานพร้อมและสามารถนำมาประกอบการประเมินได้ทันที, \textbf{มีบางส่วน} หมายถึงหลักฐานมีอยู่แล้วแต่ยังไม่ครบถ้วนหรือยังต้องการการยืนยันเพิ่มเติม และ \textbf{MARK} หมายถึงหลักฐานยังไม่มีหรือยังไม่ครบ จำเป็นต้องดำเนินการจัดหาหรือจัดทำเพิ่มเติมก่อนการประเมินภายนอก

\section{สรุปจำนวนหลักฐานตามสถานะ}

\noindent\begin{tabularx}{\linewidth}{|>{\raggedright\arraybackslash}X|c|c|}
\hline
\textbf{สถานะ} & \textbf{จำนวน (รายการ)} & \textbf{ร้อยละ} \\
\hline
มี & 20 & 51\% \\
มีบางส่วน & 14 & 36\% \\
MARK (ยังไม่มี) & 5 & 13\% \\
\hline
\textbf{รวม} & \textbf{39} & \textbf{100\%} \\
\hline
\end{tabularx}

\bigskip

\textbf{หมายเหตุ:} หลักฐานที่มีสถานะ MARK จำนวน 5 รายการ (ร้อยละ 13) เป็นหลักฐานที่ต้องดำเนินการจัดหาหรือจัดทำเพิ่มเติม โดยเฉพาะอย่างยิ่ง มคอ.3 รายวิชาบังคับ แผนการสอน 6 วิชา และแผน Alumni Tracking — หลักฐาน Criterion 8 เพิ่มเติม (AUNQA-8-4, 8-5) อยู่ระหว่างสะสมผล CLO/ความพึงพอใจ

\bigskip

\section{ตารางสรุปความต้องการหลักฐานตาม Criterion}

\setlength{\tabcolsep}{3pt}
\begin{tabularx}{\linewidth}{|>{\raggedright\arraybackslash}p{3.2cm}|Y|>{\raggedright\arraybackslash}p{3.5cm}|}
\hline
\textbf{Criterion} & \textbf{หลักฐานหลักที่ต้องการ} & \textbf{สถานะ} \\
\hline
Criterion 1 & AUNQA-1-1 -- 1-4 (มคอ.2, ประชุมหลักสูตร, CLO--PLO, มติสภา) & \cellcolor{green!25}มี \\
Criterion 2 & AUNQA-2-1 -- 2-4 (มคอ.2, CLO--PLO, ข้อบังคับบัณฑิตศึกษา, มคอ.3 บังคับ) & \cellcolor{yellow!25}มีบางส่วน \\
Criterion 3 & AUNQA-3-1 -- 3-5 (มคอ.2 หมวด 4--5, CLO 32 วิชา, มคอ.3, ข้อบังคับ, แผนสอน 6 วิชา) & \cellcolor{yellow!25}มีบางส่วน \\
Criterion 4 & AUNQA-4-1 -- 4-3 (ประกาศอุทธรณ์, ข้อบังคับบัณฑิตศึกษา, ประกาศส่งผลการประเมิน) & \cellcolor{green!25}มี \\
Criterion 5 & AUNQA-5-1 -- 5-4 (แผนบุคคล, นโยบายพัฒนาอาจารย์, ข้อบังคับบุคคล, รายงานพัฒนาอาจารย์) & \cellcolor{yellow!25}มีบางส่วน \\
Criterion 6 & AUNQA-6-1 -- 6-5 (ข้อบังคับรับเข้า, ประกาศรับสมัคร, ห้องสมุด, ความพึงพอใจ, กิจกรรมเสริม) & \cellcolor{yellow!25}มีบางส่วน \\
Criterion 7 & AUNQA-7-1 -- 7-5 (รายงานสิ่งอำนวยความสะดวก, ฐานข้อมูล, แผน IT, ประเมินคุณภาพ, บันทึกห้องปฏิบัติการ) & \cellcolor{yellow!25}มีบางส่วน \\
Criterion 8 & AUNQA-8-1 -- 8-5 (REG, Alumni, ผลงานอาจารย์, สรุป CLO/PLO ระหว่างทาง, แบบสำรวจความพึงพอใจ) & \cellcolor{yellow!25}มีบางส่วน \\
\hline
\end{tabularx}

\section*{หมายเหตุสิ้นสุดรายงาน}
\addcontentsline{toc}{section}{หมายเหตุสิ้นสุดรายงาน}

\bigskip

รายงานประเมินตนเองตามเกณฑ์ AUN-QA ฉบับนี้จัดทำขึ้นเป็น \textbf{Initial Self-Assessment} โดยหลักสูตรวิศวกรรมศาสตรมหาบัณฑิต สาขาวิชาวิศวกรรมคอมพิวเตอร์และปัญญาประดิษฐ์ (หลักสูตรใหม่ พ.ศ. 2568) เริ่มรับนักศึกษาในปีการศึกษา 2568 จำนวน 4 คน (เข้าศึกษาภาคการศึกษาที่ 2) คะแนนประเมินตนเองส่วนใหญ่จึงยังเป็นระดับ 2 (มีระบบและการดำเนินงาน กำลังสะสมผลจากการจัดการเรียนการสอน) ยกเว้น Criterion 8 ที่ได้ระดับ 1 เนื่องจากยังไม่มีบัณฑิตและยังไม่มีผลลัพธ์ระดับหลักสูตรที่สรุปเป็นตัวเลขได้ครบ

คณะผู้รับผิดชอบหลักสูตรมุ่งมั่นที่จะพัฒนาคุณภาพอย่างต่อเนื่องและเตรียมความพร้อมให้ครบถ้วนตามเกณฑ์ AUN-QA ก่อนการประเมินภายนอกโดยคณะกรรมการจากสภามหาวิทยาลัยและ/หรือองค์กรภายนอก

\bigskip

\begin{flushleft}
\textbf{ลงชื่อ}\\[4pt]
................................................\\[2pt]
ประธานหลักสูตรวิศวกรรมคอมพิวเตอร์และปัญญาประดิษฐ์\\
คณะเทคโนโลยีอุตสาหกรรม\\
มหาวิทยาลัยราชภัฏอุตรดิตถ์\\[10pt]
\textbf{วันที่} ....... / ....... / .......
\end{flushleft}
